\documentclass[11pt]{article}
\usepackage[utf8]{inputenc}
\usepackage{geometry}
\geometry{a4paper, margin=1in}
\usepackage{booktabs}
\usepackage{hyperref}

\title{Estimating the Causal Effect of Insecticide-Treated Nets Use on Malaria Infection Among Children Under Five in Kenya}
\author{Nicholas Mwanza}


\begin{document}
	\maketitle
	\today
	
	\section*{Overview}
	
	This document explains what each of my R scripts does. I analyzed data from the Kenya Malaria Indicator Surveys (KMIS) from 2015 and 2020 to estimate the effect of insecticide-treated net (ITN) use on malaria among children under 5.
	
	The data comes from three sources: individual-level malaria test results (PR files), household-level information (HR files), and environmental data (geocov files). I merged these into one analytical dataset.
	
	Below is a brief description of each script and what it produces.
	
	\section*{Week 1: 01\_data\_loading.R}
	
	\textbf{What it does:} Loads all raw KMIS data for 2015 and 2020. Cleans the environmental data by replacing negative values (-9999) with missing values. Filters to children under 5. Merges PR, HR, KR, shapefile, and geocov data into one dataset. Also checks if the sample is large enough for Double Machine Learning.
	
	\textbf{What it produces:} \texttt{analytical\_2015.rds}, \texttt{analytical\_2020.rds}, \texttt{analytical\_combined.rds}
	
	\textbf{Sample size after this step:} 6,771 children (3,286 in 2015, 3,485 in 2020)
	
	\section*{Week 2: 02\_data\_cleaning.R}
	
	\textbf{What it does:} Takes the raw merged data and properly formats variables. Converts wealth into 5 categories (Poorest to Richest), residence into Urban/Rural, sex into Male/Female. Creates binary variables for electricity, improved water, and improved sanitation. Extracts GPS coordinates from shapefiles for mapping. Creates cluster-level datasets (one row per cluster with malaria prevalence and ITN use rate).
	
	\textbf{What it produces:} Cleaned child-level datasets, cluster-level datasets for mapping, and shapefiles with environmental data merged.
	
	\section*{Week 3: 03\_descriptive\_analysis.R}
	
	\textbf{What it does:} Generates summary statistics using survey weights. Creates tables and figures showing malaria prevalence and ITN use by wealth, residence, and age. Makes maps of malaria prevalence and ITN use across Kenya.
	
	\textbf{Key findings:} Malaria prevalence dropped from 8.2\% in 2015 to 4.8\% in 2020. ITN use increased from 54.2\% to 57.1\%. Malaria is highest in poorer households and among children aged 12-23 months.
	
	\textbf{What it produces:} Descriptive tables, line plots, bar charts, and spatial maps (PNG files).
	
	\section*{Week 4: 04\_baseline\_regression.R}
	
	\textbf{What it does:} Runs traditional logistic regression, GLMM (mixed models), and GLMM-based causal methods (Propensity Score Matching, IPTW, Doubly Robust). This is the most detailed script because I had to address 13 supervisor comments.
	
	\textbf{Key findings:} GLMM shows 46-49\% reduction in malaria odds with ITN use. PSM gives around 54\% reduction. The high ICC (41-47\%) justifies using multilevel models.
	
	\textbf{What it produces:} Regression tables, caterpillar plots for random effects, love plots for balance checks, overlap plots for propensity scores, and a master summary table combining all methods.
	
	\section*{Week 5: 05\_double\_ml.R}
	
	\textbf{What it does:} Implements Double Machine Learning with cluster-aware cross-fitting (keeps clusters together during splitting). Reports ATE (risk difference) instead of odds ratios. Calculates E-values for sensitivity analysis (VanderWeele \& Ding 2017). Uses Random Forest and XGBoost as ML learners.
	
	\textbf{Key findings:} DML estimates a population-average effect of 2-4 percentage point reduction in malaria risk (OR around 0.93-0.97). This is smaller than GLMM conditional effects, which is expected given the high clustering. E-values around 1.2-1.3 suggest modest robustness to unmeasured confounding.
	
	\textbf{What it produces:} DML results table with ATE, standard errors, p-values, and E-values. Forest plot on ATE scale. E-value bar plot.
	
	\section*{Week 6: 06\_robustness\_heterogeneity.R}
	
	\textbf{What it does:} Tests whether DML estimates change with different ML learners (tuned RF, smaller node size, XGBoost with more rounds). Also examines effect modification by wealth quintile, urban/rural residence, age group, and rainfall level.
	
	\textbf{Key findings:} Estimates are robust across learners. Protective effect is strongest in poorer households and among children aged 12-23 months. Weaker effect in richest households (almost null). No meaningful difference by rainfall level.
	
	\textbf{What it produces:} Robustness comparison table and plot. Heterogeneity forest plot showing ORs for all subgroups.
	
	\section*{Summary of Results}
	
	\begin{table}[htbp]
		\centering
		\begin{tabular}{ll}
			\toprule
			\textbf{Method} & \textbf{Effect (OR range)} \\
			\midrule
			GLMM (conditional) & 0.51 - 0.54 \\
			GLMM-based PSM & 0.46 \\
			GLMM-based IPTW/DR & 0.50 - 0.53 \\
			DML (population-average) & 0.93 - 0.98 \\
			\bottomrule
		\end{tabular}
		\caption{Comparison across methods}
	\end{table}
	
	ITN use reduces malaria risk. The effect is stronger when accounting for cluster-level variation (conditional effects) and weaker at the population level (ATE), which is expected given the high ICC.

\end{document}